\section{Moment exactness: the proof route, the conjecture and the shot bound}
\label{sm:moments}

That proof is three lines of Rayleigh--Ritz, and it is not the proof this statement is usually given.
The usual route argues that an order-$(K{+}1)$ quadrature ``matches the first $2K+1$ moments
exactly'', and $2K+1$ is exactly what Table~\ref{tab:moments} measures: the \emph{reach} is right,
the route is not. What does not hold here
is the identification of $A_{\mathcal S}$ with a Gaussian-quadrature rule built by a nonorthogonal
Lanczos/Krylov construction. No such object is built anywhere in this work: every reconstruction
diagonalizes $P_{\mathcal S}HP_{\mathcal S}$ on a coordinate subspace, so $d\mu_{\mathcal S}$ carries
$|\mathcal S|$ nodes, not $K+1$, and the orders that separate the two objects measure the distance between them,
not an error in the published reach. The conclusion is right, the route is not, and the correct route
is also the more general one (Sec.~\ref{app:moments:exact}).

Equation~\eqref{eq:momex} invites a corollary: ``consequently, at fixed broadening
$\eta>0$, the smoothed reconstruction converges geometrically,
$\lVert A-A_{\mathcal K_K}\rVert_{1}\le C(\eta)\,\rho(\eta)^{-K}$.'' It does not follow, and it is
not a corollary here. Moment exactness is a list of $2K+2$ linear identities, and the $L_1$
distance between two Lorentzian-broadened measures is not a finite linear functional of them;
moment exactness supplies neither of the two ingredients a conversion would need
(Sec.~\ref{app:moments:conj}), and neither $C(\eta)$ nor
$\rho(\eta)$ is derived or exhibited anywhere in this work or in the deposit. We therefore state the
claim as what it is.

One constraint on any answer is already available. The two-level family of Sec.~\ref{sec:retraction}
has $\mathcal K_0=\operatorname{span}\{\phi\}$, and by direct integration at $\eta=0.15$ its
$\lVert A-A_{\mathcal K_0}\rVert_{1}$ is $1.983$ at $g/\eta=2\times10^{2}$ and $1.998$ at
$g/\eta=6.7\times10^{3}$, rising to $2\lVert\phi\rVert^{2}$. Evaluating Conjecture~\ref{conj:geom} at
$K=0$ therefore forces $C(\eta)\ge2\lVert\phi\rVert^{2}$: the same family that forces the trivial
constant on the weight-only bound forces it here on the prefactor. Whatever content the conjecture has
must be carried by $\rho(\eta)^{-K}$ alone, and in particular it cannot be non-vacuous at the small $K$
at which a shot-limited subspace operates.

One figure of this work is read against that conjecture and not against a theorem. A geometric collapse $\text{rel-}L_1\sim\rho^{-K}$ with $\rho\approx1.7$, fitted to panel~(b) of Fig.~\ref{fig:method}, is easily reported as ``precisely the geometric convergence guaranteed by the
moment-exactness theorem''. It is guaranteed by no such thing---that is the corollary ruled out
above---and the fit does not carry the claim on its own either. The panel is one of the two places in this work where
the sampling is genuine and its budget is declared---multinomial draws at $4000$ shots per time step,
eight seeds, $\eta=0.15\,t$, so $T=4000(K{+}1)$ and $T=5.2\times10^{4}$ at the last
point---and the fit uses six of its thirteen measured points, inside a window selected after the fact
($4\le K\le9$, set as a literal in \texttt{spectral\_validation\_max.py}). Across that window the mean
subspace the sampler returns grows from $179$ to $236$ determinants of the $300$-dimensional sector, so
decay in the Krylov order is confounded with growth in $|\mathcal S|$; and the $K$ of the abscissa is
the number of measured time slices, not the $K_{\mathcal S}$ of the weight-only bound, which is $-1$ at
all thirteen points by Eq.~\eqref{eq:exclusion}. No rate is fitted to that window here, and the
caption claims none.
What the panel does show, over all thirteen points, with its eight-seed band and with $T$ attached, is
that the error falls by a factor $32$ as the measured Krylov order is climbed at fixed $\eta$, from
$0.59$ at $K=0$ to $0.0185$ at $K=12$, with a single non-monotone step at $K=3$ inside the discovery
phase; that is the statement its caption makes.

There is a sharper reason not to read Eq.~\eqref{eq:momex} as an error bound, and
Table~\ref{tab:moments} contains it. For a coordinate subspace $P_{\mathcal S}$ is diagonal, so
$\operatorname{span}(\mathcal S)\supseteq\mathcal K_K$ at any $K\ge0$ already forces
$\operatorname{supp}(\phi)\subseteq\mathcal S$, that is $w_{\mathcal S}=1$ exactly---as the fourth
column records for all three rows. Moment exactness here is strictly stronger than unit captured
weight, and unit captured weight is precisely the condition under which Sec.~\ref{sec:leakage} measures
the error to be $O(1)$. The two meet in the first row: the $K=0$ member of the ladder \emph{is} the
witness $\mathcal S^{\ast}=\operatorname{supp}(\phi)$ of Sec.~\ref{sec:leakage}, of size $(\tfrac12+\tfrac1L)D$ at $L=6$ and fixed by the probe rather than by any ranking. It is moment-exact through
order $1$, it holds the entire Born weight of the probe, and at $\eta=0.18\,t$ its relative $L_1$ error
is $0.28$ on the open chain of Table~\ref{tab:moments} at $L=6$, and $0.43$ at $L=6$ and $0.35$ at $L=8$ on the periodic rings on which Sec.~\ref{sec:leakage} reports it: a measured range of
$0.28$--$0.43$ at unit captured weight, whose lower end is the very system on which the method
validation of Sec.~\ref{sec:method} is performed. (All three are $\lVert A-A_{\mathcal S}\rVert_1$
divided by $\lVert\phi\rVert^{2}$, the normalization of Theorem~\ref{thm:leakage}; the size of the
conversion to the repository's own in-window relative $L_1$ is measured in
Sec.~\ref{sec:retraction}, and the two conventions are never interchanged silently.) Exactness of $2K+2$ numbers is
what the structure certifies; the curve is certified, at fixed $\eta$ and on the subspace in hand, by
Theorem~\ref{thm:leakage}, whose own non-vacuity is measured rather than assumed
(Sec.~\ref{sec:frontier}).

What Eq.~\eqref{eq:shots} does \emph{not} say belongs in the same breath, and in the main text rather
than in an appendix. The parameter $\varepsilon$ thresholds a Born probability; it is not a spectral
accuracy, and no step of the argument connects it to $\lVert A-A_{\mathcal S}\rVert_{1}$. Arbitrarily
many configurations may each carry $p(x)<\varepsilon$ while jointly carrying most of the mixture; and,
decisively, even at $w_{\mathcal S}=1$ exactly---nothing left to capture at any $\varepsilon$---the
relative error is $O(1)$, as Sec.~\ref{sec:leakage} measures and the first row of
Table~\ref{tab:moments} exhibits. Capturing the high-probability configurations is a statement about
the sampler, not about the spectrum.

Two compositions of the bound follow easily from it, and neither is made here. ``A sampling-capture
bound polynomial in $|\mathcal S|$ and independent of the Hilbert-space dimension'' is a pair of
clauses that are literally true and jointly empty, because this same work
measures $|\mathcal S|\sim e^{1.05L}$ for these lattices, Eq.~\eqref{eq:supportexp}, so a budget
polynomial in $|\mathcal S|$ is exponential in $L$, and independence of the Hilbert-space dimension is
no saving when the quantity the budget is polynomial in grows exponentially itself. And the fall of the
required sector fraction with $L$ is not ``the favourable face of the
polynomial-shot bound'': it is not a face of that bound at all. Equation~\eqref{eq:shots} relates a
shot budget to a Born-probability threshold and contains no statement about $|\mathcal S|/D$, about
$L$, or about any trend in either.

What stands in their place is the narrower statement Eq.~\eqref{eq:shots} does support, and it is the stronger
one for being narrow: that many shots buy, at confidence $1-\delta$, every configuration above Born
weight $\varepsilon$, and nothing further. What the resulting subspace costs, and what it is worth, are
separate measurements, reported as three independent prices---support, resolution and shots---in
Sec.~\ref{sec:prices}.
